\originalchapter{22.02 作业八:演化与伽马级联}
题源: Spring 2012 Problem Set 8, 2012年4月23日. 题 4、5 原 PDF 含示范答案, 其余补写. 用原符号 n 标记径向本征态序号.
\begin{exercise}\label{ex:2202-8-1}
原题 1. 球无限阱 $V(r<R_0)=0$, 外部无限. 首行初态为 $|1,0,0\rangle/\sqrt2+\e^{\ii\pi/3}|1,1,-1\rangle/2+|2,1,0\rangle/2$, 但紧接的展开式第三项写 $R_{2,0}Y_{0,0}/2$. (a) 测 $L^2=2\hbar^2,6\hbar^2$ 的概率? (b) $\langle L_z(t)\rangle$? (c) 改初态为 $|1,0,0\rangle/\sqrt2+\e^{\ii\pi/3}|3,0,0\rangle/2+|2,0,0\rangle/2$, 求演化及半径在 $R_0/2$ 到 $R_0/2+\dd r$ 的概率. (d) 测 $E_{1,0},E_{2,1},E_{3,0}$ 的概率和平均能演化? (e) 平均半径演化?
\end{exercise}
\begin{solution}
(a) 原题内部两式不一致: 按 ket 行第三项 l=1, 则 P(l=1)=1/2, P(l=2)=0; 按展开式第三项 l=0, 则 P(l=1)=1/4, P(l=2)=0. 两者均归一, 不能擅自声称只有一个答案. (b) 两版本 $\langle L_z\rangle=-\hbar/4$, 因只有第二项 m=-1, 与时间无关.

(c) 新态三项均 l=0, 无该歧义. 归一径向 $u_n=\sqrt{2/R_0}\sin(n\pi r/R_0)$, $E_n=n^2E_1$, $E_1=\hbar^2\pi^2/(2mR_0^2)$. 记 $\omega=E_1/\hbar$, 系数 $c_1=1/\sqrt2,c_2=1/2,c_3=\e^{\ii\pi/3}/2$, 则 $\psi=\sum_{n=1}^3c_nu_nY_{00}\e^{-\ii n^2\omega t}/r$. 全角积分后径向概率密度 $P_r=|\sum_nc_nu_n\e^{-\ii n^2\omega t}|^2$. 在 $r=R_0/2$, $u_2=0,u_3=-u_1$, 因而
\begin{equation}
 P_r(R_0/2,t)\dd r=\frac{2\dd r}{R_0}
 \left[\frac34-\frac1{\sqrt2}\cos(8\omega t-\pi/3)\right].
\end{equation}
这是球壳概率, 不是某个角方向的小体积概率, 后者还要 $|Y_{00}|^2/r^2$.

(d) 对 (c) 新态概率为 $1/2,0,1/4$, 另 $E_{2,0}$ 为 $1/4$. 平均能 $\langle E\rangle=(1/2+4/4+9/4)E_1=15E_1/4$, 时间不变. 若题意仍对最初态问 (d), 按 ket 版 $E_{1,0},E_{2,1}$ 概率 $1/2,1/4$, $E_{3,0}$ 为0, 另 $E_{1,1}$ 为 $1/4$; 本稿分别明示依赖初态.

(e) 对 (c), 对角半径矩阵元全为 $R_0/2$, 不同 n,m 时
\begin{equation}
 r_{nm}=\frac{R_0}{\pi^2}\left[\frac{(-1)^{n-m}-1}{(n-m)^2}-\frac{(-1)^{n+m}-1}{(n+m)^2}\right].
\end{equation}
所以 $r_{12}=-16R_0/(9\pi^2)$, $r_{23}=-48R_0/(25\pi^2)$, $r_{13}=0$. 平均半径为
$R_0/2+r_{12}\cos(3\omega t)/\sqrt2+r_{23}\cos(5\omega t-\pi/3)/2$. 概率干涉随时间变化但能量权重恒定.
\end{solution}
\begin{exercise}\label{ex:2202-8-2}
原题 2. 给全部允许 gamma 多极及可能主导者: (a) $5/2^-\to7/2^+$, (b) $4^+\to2^+$, (c) $9/2^+\to3/2^-$, (d) $3/2^-\to7/2^-$, (e) $3^+\to2^-$, (f) $2^+\to2^+$.
\end{exercise}
\begin{solution}
按 L 范围与宇称, (a) E1,M2,E3,M4,E5,M6, 通常 E1 主导; (b) E2,M3,E4,M5,E6, 通常 E2; (c) E3,M4,E5,M6, 通常 E3; (d) E2,M3,E4,M5, 通常 E2; (e) E1,M2,E3,M4,E5, 通常 E1; (f) M1,E2,M3,E4, 通常 M1 或 E2, 取决于核矩阵元. ``最低 L 较强''是长波估算, 遇禁戒或组态抑制可改变, 不能把选择规则当精确强度计算.
\end{solution}
\begin{exercise}\label{ex:2202-8-3}
原题 3. $^{81}$Br 级联: 第二激发态到第一为260 keV M2, 第一到基态276 keV M1、E2. 价质子在28--50壳的四轨道中, 结合选择规则和配对确定三态自旋宇称.
\end{exercise}
\begin{solution}
四轨道为 $p_{3/2},f_{5/2},p_{1/2},g_{9/2}$. M1/E2 不变宇称, 所以前两态均负宇称; M2 改变宇称, 所以第二激发是唯一正宇称 $9/2^+$. 它经 L=2 必须到 $j\ge5/2$, 负宇称三轨道中只有 $5/2^-$ 能承接. 第一激发故为 $5/2^-$, 基态由课程壳次序和配对为 $3/2^-$. 能量分别0,276,536 keV. 单凭选择规则不能把基态 $1/2^-$ 完全排除, 采用课程价壳信息是必要的附加条件.
\end{solution}
\begin{exercise}\label{ex:2202-8-4}
原题 4 (含答案). $^{211}$Po 经 alpha 留 $^{207}$Pb. (a) 单纯填壳预测自旋宇称? (b) 配对重排后的基态? (c) $3/2^-$ 的897 keV态到基态允许哪些多极, 两主要率之比? 同时比较到 naive $13/2^+$ 情形. (d) 观察328 keV M1与569 keV E2, 求中间态及能量. (e) 画未闭壳占据示意说明基态与两激发态.
\end{exercise}
\begin{solution}
官方答案重编. (a) 填壳示意最后奇核子在 $i_{13/2}$, 给 $13/2^+$. (b) 126中子闭壳附近允许配对重排, 低 j 的 $p_{1/2}$ 空穴给 $1/2^-$. (c) $3/2^-\to1/2^-$ 允许 M1、E2. 源估计 $\lambda(E2)\simeq7.3\times10^7A^{4/3}E_\gamma^5$, $\lambda(M1)\simeq5.6\times10^{13}E_\gamma^3$ s$^{-1}$, 能量用 MeV, 给比约 $1.3\times10^{-6}A^{4/3}E_\gamma^2\simeq1.3\times10^{-3}$, M1 较强. 到 $13/2^+$ 时 L=5--8, 允许 E5,M6,E7,M8, 低阶 E5 通常主导但极受抑.

(d) 中间态 $5/2^-$, 能量约569 keV, 级联图见第 \ref{ch:beta-gamma} 章. (e) 以126闭壳为参照, 基态在 $3p_{1/2}$ 留一个中子空穴, 第一激发在 $2f_{5/2}$ 留一空穴, 第二在 $3p_{3/2}$ 留一空穴, 其余近闭壳轨道成对. 可编辑的空穴图如下, 圆圈指空穴, 不是额外实粒子:
\begin{center}\begin{tikzpicture}[x=1cm,y=0.7cm]
\foreach\x/\t in{0/基态,3/第一激发,6/第二激发}{\node at(\x+0.8,3.2){\t};\foreach\y/\lab in{0/p_{3/2},1/f_{5/2},2/p_{1/2}}{\draw(\x,\y)--(\x+1.5,\y);\node[left]at(\x,\y){$\lab$};}}
\draw(0.8,2)circle(0.10);\draw(3.8,1)circle(0.10);\draw(6.8,0)circle(0.10);
\end{tikzpicture}\end{center}
垂直位置只排组态标签, 不表示这三轨道的精确单粒子能量次序.
\end{solution}
\begin{exercise}\label{ex:2202-8-5}
原题 5 (已解). 自由粒子初态高斯 $\psi=(2\pi a^2)^{-1/4}\e^{-x^2/(4a^2)}$. (a) H 的本征态、本征值; (b) 展开初态; (c) 推 $\psi(x,t)$; (d) 给位置概率宽度演化并画初始与后来波包.
\end{exercise}
\begin{solution}
官方示范重新归一. (a) 平面波 $\e^{\ii kx}/\sqrt{2\pi}$, $E_k=\hbar^2k^2/(2m)$. (b) Fourier 系数 $(2a^2/\pi)^{1/4}\e^{-a^2k^2}$. (c) 逐分量乘 $\e^{-\ii E_kt/\hbar}$ 并做复高斯积分, 得第 \ref{ch:golden} 章公式. (d) 方差 $\alpha(t)^2=a^2+\hbar^2t^2/(4m^2a^2)$. 以 $a=1$ 无量纲示意初始标准差1与后来2, 两曲线面积都1:
\begin{center}\begin{tikzpicture}\begin{axis}[width=8cm,height=4.5cm,xlabel={$x/a$},ylabel={归一概率密度},domain=-6:6,samples=100]
\addplot[thick,structurecolor]{exp(-x^2/2)/sqrt(2*pi)};\addplot[thick,dashed]{exp(-x^2/8)/sqrt(8*pi)};\end{axis}\end{tikzpicture}\end{center}
宽度增加同时峰降低, 不会使总概率增加.
\end{solution}
